\documentclass{article}

% ============================================================
% NeurIPS-style formatting (self-contained, no external .sty
% dependency — mirrors the official neurips_2024 layout:
% single-column, Times, numbered sections, abstract, checklist).
% ============================================================
\usepackage[margin=1in]{geometry}
\usepackage{times}
\usepackage{amsmath,amssymb}
\usepackage{graphicx}
\usepackage{booktabs}
\usepackage{algorithm}
\usepackage{algpseudocode}
\usepackage{hyperref}
\usepackage{xcolor}
\usepackage{caption}
\usepackage[numbers,sort]{natbib}

\hypersetup{colorlinks=true, linkcolor=blue, citecolor=blue, urlcolor=blue}

\title{Persistent Kernel Scheduling for Long-Running Agentic Inference Loops}

\author{
  Manav Sutar \\
  Single Core Labs \\
  Pune, India \\
  \texttt{manav@singlecorelabs.com}
}

\begin{document}
\maketitle

\begin{abstract}
Autonomous LLM agents execute long, iterative loops of the form
\emph{model call} $\rightarrow$ \emph{tool call} $\rightarrow$ \emph{model call},
often for tens to hundreds of steps per task. Standard inference stacks
re-launch a full stack of CUDA kernels at every model call, paying
launch and framework dispatch overhead on every step even though the
underlying transformer weights and much of the on-GPU state persist
across the loop. Recent \emph{megakernel} and \emph{persistent kernel}
techniques eliminate per-operator launch overhead for single-turn
decoding by fusing an entire forward pass into one resident kernel
that is fed by an in-kernel task queue \citep{mpk2025,kog2026,adamk2026}.
However, these designs target token-by-token decoding within a single
generation; they do not address the distinct overhead structure of
\emph{agent loops}, where the GPU sits idle across externally-latent
tool calls and where control flow depends on tool output rather than
token output. We formalize this gap, propose a persistent-kernel
scheduling design that keeps a resident kernel warm across tool-call
boundaries via a host-managed work queue, and provide a reference
implementation together with a benchmarking harness that isolates
launch-overhead cost specifically in multi-step tool-calling loops
(as opposed to raw decode). We report a CPU-level simulation
methodology for estimating overhead scaling with step count and
tool latency, and provide CUDA/Triton skeleton kernels intended to
be benchmarked on target hardware (evaluated here on consumer-class
GPUs; A100/H100-class results are left as an open experiment for
readers with access to that hardware).
\end{abstract}

\section{Introduction}

Two workloads are routinely conflated in the megakernel literature:
(1) single-turn autoregressive decoding, where the loop body is
entirely on-GPU and iteration count is bounded by output length, and
(2) agentic loops, where each iteration alternates an on-GPU decode
segment with an off-GPU, data-dependent tool call (a database query,
an API call, a subprocess). Prior megakernel systems --- Mirage
Persistent Kernel (MPK) \citep{mpk2025}, the AMD MI300X monokernel
\citep{kog2026}, and Ada-MK's DAG-search fusion \citep{adamk2026} ---
report that kernel-launch overhead consumes roughly 5--15\% of
end-to-end decode latency, and that fusing the decode forward pass
into one resident kernel closes most of that gap. Event Tensor
\citep{eventtensor2026} further shows that the harder open problem is
\emph{data-dependent} dynamism inside a megakernel, i.e., control flow
whose shape is not known until runtime.

Agent loops push this data-dependence outside the kernel entirely: the
next model call's *content* depends on a tool result the GPU cannot
see until the host receives it. This raises three questions distinct
from single-turn decode optimization:

\begin{enumerate}
  \item Can a persistent kernel remain resident (registers/shared
  memory allocated, warps parked) across a tool call whose latency is
  unknown at compile time, rather than being torn down and relaunched?
  \item What is the actual overhead contributed by kernel
  \emph{teardown/relaunch} versus unavoidable tool latency, as a
  function of step count $T$ and tool latency $L$?
  \item Can multiple concurrent agent instances share one resident
  kernel via warp specialization, batching their decode segments while
  their tool calls are in flight independently?
\end{enumerate}

This paper's contribution is threefold: (a) a cost model separating
launch overhead from tool latency in agent loops (\S\ref{sec:method});
(b) a persistent-kernel design with a host-managed work queue that
keeps the kernel warm across tool boundaries (\S\ref{sec:design}); and
(c) a reference implementation and benchmarking harness
(\S\ref{sec:impl}), with results left parameterized for the reader's
own hardware (\S\ref{sec:experiments}).

\section{Related Work}

\textbf{Kernel fusion for decode.} MPK \citep{mpk2025} compiles
multi-GPU inference into an SM-level task graph executed inside a
single mega-kernel, reporting 1.0--1.7$\times$ throughput gains over
kernel-per-operator systems such as vLLM and SGLang. The Kog AI
monokernel \citep{kog2026} implements an entire decode pass as one
persistent kernel on AMD MI300X, reporting 3{,}000+ tokens/s per
request. Ada-MK \citep{adamk2026} profiles kernel-launch overhead at
roughly 14.6\% of end-to-end latency for a 1.5B model under
TensorRT-LLM and proposes automated DAG search to decide fusion
boundaries.

\textbf{Dynamic megakernels.} Event Tensor \citep{eventtensor2026}
introduces a compiler abstraction encoding data-dependent task
dependencies, addressing the case where tile shapes are not known
statically --- the closest existing treatment of the dynamism problem
that agent loops exhibit at a coarser (tool-call) granularity.

\textbf{GPU-resident runtimes.} GPUOS \citep{gpuos2026} proposes a
long-lived kernel fed by a host-managed work queue with runtime
operator injection via NVRTC, avoiding kernel relaunch for
heterogeneous, evolving workloads. Our design borrows this
host-queue/resident-kernel split but targets the agent-loop boundary
specifically, i.e., the queue entries are \emph{tool-call results},
not arbitrary injected operators.

\textbf{Gap.} No existing system we are aware of measures or
optimizes for the specific overhead of kernel relaunch \emph{across
tool-call boundaries} in agent loops, as opposed to across
token-generation steps within one decode. This is the gap this paper
targets.

\section{Cost Model}
\label{sec:method}

Let an agent loop consist of $T$ steps. At step $t$, the model
performs a decode segment of on-GPU latency $d_t$, followed by a tool
call of latency $\ell_t$ (network- or CPU-bound, effectively opaque to
the GPU). Under a naive per-step kernel-launch design, each decode
segment additionally pays a launch/teardown overhead $\kappa$, so
total wall-clock time is

\begin{equation}
  T_{\text{naive}} = \sum_{t=1}^{T} \left( d_t + \kappa + \ell_t \right).
  \label{eq:naive}
\end{equation}

Under a persistent-kernel design where the kernel remains resident
across tool calls (parked on an idle work-queue poll rather than
torn down), the per-step overhead collapses to a lightweight queue
signal $\sigma \ll \kappa$:

\begin{equation}
  T_{\text{persistent}} = \sum_{t=1}^{T} \left( d_t + \sigma + \ell_t \right).
  \label{eq:persistent}
\end{equation}

The recoverable overhead is therefore

\begin{equation}
  \Delta(T) = T \cdot (\kappa - \sigma).
  \label{eq:delta}
\end{equation}

Two regimes matter in practice:

\begin{itemize}
  \item \textbf{Tool-latency-dominated} ($\ell_t \gg \kappa$): typical
  for agents calling slow external APIs. Here $\Delta(T)$ is a small
  fraction of wall-clock time and persistent kernels offer limited
  end-to-end speedup, though they still free the GPU for other work
  during the idle interval if the queue design allows preemption.
  \item \textbf{Launch-overhead-dominated} ($\kappa \gtrsim d_t$):
  typical for short tool calls (cache lookups, local function calls,
  small retrievals) chained across many steps, or for small/quantized
  models where $d_t$ itself is only slightly larger than $\kappa$.
  Here $\Delta(T)$ scales linearly with $T$ and dominates end-to-end
  latency for long agent trajectories ($T$ in the tens to hundreds).
\end{itemize}

This distinguishes the agent-loop regime from single-turn decode: $T$
is bounded by \emph{agent steps}, not \emph{output tokens}, and
$\ell_t$ (tool latency) is a new term absent from prior megakernel
cost models.

\section{Persistent Kernel Design for Agent Loops}
\label{sec:design}

We propose a three-part design:

\textbf{(1) Resident decode kernel.} A single persistent kernel
(following the MPK/Kog pattern) implements the model forward pass as
an SM-level task graph, launched once per agent session rather than
once per step.

\textbf{(2) Host-managed agent work queue.} Between decode segments,
the kernel does not exit; warps poll a lightweight, host-writable
queue (cf. GPUOS \citep{gpuos2026}) for the next input token block,
which becomes available only once the host has received the tool
result and re-encoded it. This converts step boundaries from
"relaunch" events into "queue poll" events.

\textbf{(3) Cooperative yield on tool latency.} Because $\ell_t$ can
be large and unpredictable, the resident kernel exposes a yield path:
if no queue entry arrives within a short poll window, the relevant
SMs are released back to the scheduler (e.g., via MPS or
time-slicing) so concurrent agent sessions --- or unrelated
workloads --- can use the GPU during the idle interval, then
re-acquire the resident context when the tool result arrives. This
addresses the multi-agent batching question raised in
\S\ref{sec:method}: several agent sessions can share the GPU's SMs
opportunistically without each paying a full relaunch.

Algorithm~\ref{alg:loop} summarizes the host/device split.

\begin{algorithm}[t]
\caption{Persistent-kernel agent loop}
\label{alg:loop}
\begin{algorithmic}[1]
\State \textbf{Device:} launch resident kernel $K$ once; $K$ polls queue $Q$
\For{$t = 1, \dots, T$}
  \State \textbf{Host:} push encoded state onto $Q$
  \State \textbf{Device:} $K$ dequeues, executes decode segment $d_t$, writes output tokens
  \State \textbf{Host:} read output tokens; dispatch tool call $\ell_t$ (async)
  \If{no new entry in $Q$ within poll window $w$}
    \State \textbf{Device:} yield SMs (cooperative release)
  \EndIf
  \State \textbf{Host:} on tool completion, re-encode result, go to line 3
\EndFor
\State \textbf{Device:} $K$ exits on host-issued shutdown signal
\end{algorithmic}
\end{algorithm}

\section{Reference Implementation}
\label{sec:impl}

We provide three artifacts alongside this paper (see supplementary
code):

\begin{itemize}
  \item \texttt{persistent\_kernel.cu} --- a CUDA skeleton implementing
  a resident kernel with a device-side spin-queue, a decode-segment
  stub, and a host-side loop that pushes/pops queue entries instead of
  relaunching the kernel. Intended as a scaffold to plug in a real
  forward pass (e.g., via CUTLASS or a Mirage-compiled task graph).
  \item \texttt{triton\_persistent\_agent\_kernel.py} --- a Triton
  prototype of the same idea at a higher level, useful for rapid
  iteration before dropping to raw CUDA.
  \item \texttt{benchmark.py} --- a hardware-agnostic benchmarking
  harness that measures $T_{\text{naive}}$ vs.\ $T_{\text{persistent}}$
  under Eq.~\ref{eq:naive}--\ref{eq:delta}. On machines without a
  CUDA-capable GPU (as in this preparation environment) it runs in a
  \emph{simulation mode} that models $\kappa$, $\sigma$, $d_t$, and
  $\ell_t$ as configurable random variables to validate the cost model
  and produce the scaling plots in \S\ref{sec:experiments}; on a real
  GPU it switches automatically to measuring actual kernel-launch and
  queue-poll latency via CUDA events.
\end{itemize}

\section{Experiments}
\label{sec:experiments}

\textbf{Setup.} We evaluate the cost model of \S\ref{sec:method}
under two parameter regimes matched to the two cases discussed above:
(A) tool-latency-dominated, $\ell_t \sim \mathcal{U}(50,200)$ ms,
$d_t \sim \mathcal{U}(5,15)$ ms, $\kappa \approx 1$--2 ms (typical
measured CUDA launch overhead in prior work); and (B)
launch-overhead-dominated, $\ell_t \sim \mathcal{U}(0.5,2)$ ms,
$d_t \sim \mathcal{U}(2,5)$ ms, same $\kappa$. In both regimes
$\sigma$ (queue-poll signal) is taken as $\kappa/20$, consistent with
the lightweight-signaling overheads reported for MPK's in-kernel
runtime \citep{mpk2025}.

\textbf{Metric.} We report $\Delta(T)/T_{\text{naive}}(T)$, the
fraction of end-to-end wall-clock time attributable to recoverable
launch overhead, as a function of trajectory length $T \in
\{1,10,50,100,200\}$.

\textbf{Placeholder for hardware results.} Table~\ref{tab:results}
reports simulated values from \texttt{benchmark.py}; column
``measured (GPU)'' is left for the reader to fill in by running the
CUDA/Triton artifacts on target hardware (we recommend at minimum an
RTX-class consumer GPU and, if available, an A100/H100 for direct
comparison against the MPK/Kog numbers cited in \S2).

\begin{table}[t]
\centering
\caption{Fraction of wall-clock time attributable to launch overhead,
by trajectory length $T$ and regime, with $\kappa=1.5$ ms and
$\sigma=\kappa/20$ (modeled queue-poll cost). ``Simulated'' values are
produced by \texttt{benchmark.py --mode simulate} on a CPU-only
machine (reproducible with the default seed); ``measured (GPU)'' is
intentionally left blank for the reader's own hardware run via
\texttt{benchmark.py --mode gpu}, which measures real $\kappa$ and
$\sigma$ with CUDA events instead of modeling them.}
\label{tab:results}
\begin{tabular}{lcccc}
\toprule
Regime & $T$ & Simulated $\Delta/T_{\text{naive}}$ & Measured (GPU) \\
\midrule
A: tool-latency-dominated    & 10  & 0.89\% & \emph{run \texttt{--mode gpu}} \\
A: tool-latency-dominated    & 100 & 1.06\% & \emph{run \texttt{--mode gpu}} \\
B: launch-overhead-dominated & 10  & 21.82\% & \emph{run \texttt{--mode gpu}} \\
B: launch-overhead-dominated & 100 & 23.18\% & \emph{run \texttt{--mode gpu}} \\
\bottomrule
\end{tabular}
\end{table}

Even in this idealized simulation, the two regimes separate cleanly:
in the tool-latency-dominated regime (A), recoverable launch overhead
is under 1.1\% of wall-clock time even at $T=100$ steps, because each
step is dominated by an opaque, un-recoverable tool-call latency
($\ell_t \in [50,200]$ ms). In the launch-overhead-dominated regime
(B) --- short tool calls chained across many steps --- the same
$\kappa$ and $\sigma$ produce a recoverable overhead of roughly
21--23\% of wall-clock time, growing slowly with $T$ toward the
asymptotic fraction $(\kappa-\sigma)/(\bar d + \kappa + \bar\ell)$.
This matches the qualitative prediction of \S\ref{sec:method}: agents
built around fast, chainable tool calls (cache lookups, local
function calls, small retrievals, as opposed to slow external-API
agents) are exactly where a persistent-kernel design is expected to
pay off, and this should be verified against real hardware
$\kappa,\sigma$ measurements via \texttt{benchmark.py --mode gpu}
rather than the modeled constants used here. We deliberately do not
fabricate GPU numbers in this draft: the harness produces them
directly when run on the target machine, and the exact values depend
on GPU generation, driver, and the forward pass plugged into the CUDA
skeleton, so hardcoding illustrative numbers into the paper would
misrepresent an actual hardware result.

\section{Discussion and Limitations}

The design in \S\ref{sec:design} addresses overhead \emph{caused by
relaunching}, not overhead caused by the tool call itself; in the
tool-latency-dominated regime the main win is freeing the GPU for
other tenants during the idle window, not reducing wall-clock latency
for a single agent. The cooperative-yield mechanism (line 6,
Algorithm~\ref{alg:loop}) is the least mature part of the design and
would benefit from integration with existing GPU scheduling
primitives (MPS, MIG, or a custom SM-partition scheme) rather than a
from-scratch implementation. We also do not address correctness under
speculative or partial decode segments that are discarded when a tool
result invalidates an in-flight assumption; this is left as future
work, connected to the broader idea of speculative multi-agent
execution.

\section{Conclusion}

Persistent/megakernel designs have closed most of the kernel-launch
overhead gap for single-turn decode. Agent loops introduce a
distinct, largely unaddressed overhead structure at the tool-call
boundary. We formalize this with a simple cost model, propose a
host-queue/resident-kernel design to close the gap, and release a
reference implementation and benchmarking harness so the cost model
can be validated empirically on real hardware rather than asserted.

\bibliographystyle{plainnat}
\begin{thebibliography}{9}

\bibitem[Mirage Persistent Kernel, 2025]{mpk2025}
Mirage Persistent Kernel: A Compiler and Runtime for Mega-Kernelizing
Tensor Programs. \url{https://arxiv.org/html/2512.22219v1}

\bibitem[Kog AI, 2026]{kog2026}
Building a single-kernel, latency-optimized LLM inference engine on
AMD MI300X GPUs. Kog AI Engineering Blog.
\url{https://blog.kog.ai/building-a-single-kernel-latency-optimized-llm-inference-engine-on-amd-mi300x-gpus/}

\bibitem[Ada-MK, 2026]{adamk2026}
Ada-MK: Adaptive MegaKernel Optimization via Automated DAG-based
Search for LLM Inference. \url{https://arxiv.org/html/2605.11581v1}

\bibitem[Event Tensor, 2026]{eventtensor2026}
Event Tensor: A Unified Abstraction for Compiling Dynamic
Megakernels. \url{https://arxiv.org/html/2604.13327v1}

\bibitem[GPUOS, 2026]{gpuos2026}
GPUOS: A GPU Operating System Primitive for Transparent Operation
Fusion. \url{https://arxiv.org/pdf/2604.17861}

\end{thebibliography}

\end{document}
